\begin{figure}[ht]
\centering
\begin{tikzpicture}[x=1cm, y=1cm, font=\small]

% Sampled points on a grid. Interior points can use the symmetric
% central stencil; the first point has no left neighbour and must use a
% one-sided (forward) stencil. The figure shows both, with the stencil
% weights labelled. No function values are computed; the picture is a
% schematic of which samples each rule reads.

% baseline with sample ticks at x = 0,1,2,3,4,5
\draw[thick] (-0.3, 0) -- (6.0, 0) node[right, font=\scriptsize] {sample index};
\foreach \i in {0,1,2,3,4,5} {
  \fill[engblack] (\i, 0) circle (1.6pt);
  \node[anchor=north, font=\scriptsize] at (\i, -0.12) {$x_{\i}$};
}

% Interior central stencil at x_3: reads x_2 and x_4 with weights -1,+1
\draw[thick, engblue, ->] (2, 0.55) -- (2, 0.12);
\draw[thick, engblue, ->] (4, 0.55) -- (4, 0.12);
\node[engblue, anchor=south, font=\tiny] at (2, 0.55) {$-\tfrac{1}{2h}$};
\node[engblue, anchor=south, font=\tiny] at (4, 0.55) {$+\tfrac{1}{2h}$};
\draw[engblue, thick, decorate, decoration={brace, amplitude=4pt}]
  (1.8, 1.05) -- (4.2, 1.05);
\node[engblue, anchor=south, font=\scriptsize] at (3, 1.12)
  {central rule at $x_{3}$ (interior)};

% Boundary one-sided (forward, 2nd order) stencil at x_0:
% reads x_0, x_1, x_2 with weights -3, +4, -1 over 2h
\draw[thick, engred, ->] (0, -1.05) -- (0, -0.15);
\draw[thick, engred, ->] (1, -1.05) -- (1, -0.15);
\draw[thick, engred, ->] (2, -1.05) -- (2, -0.15);
\node[engred, anchor=north, font=\tiny] at (0, -1.05) {$-\tfrac{3}{2h}$};
\node[engred, anchor=north, font=\tiny] at (1, -1.05) {$+\tfrac{4}{2h}$};
\node[engred, anchor=north, font=\tiny] at (2, -1.05) {$-\tfrac{1}{2h}$};
\draw[engred, thick, decorate, decoration={brace, amplitude=4pt, mirror}]
  (-0.2, -1.55) -- (2.2, -1.55);
\node[engred, anchor=north, font=\scriptsize] at (1, -1.62)
  {one-sided rule at $x_{0}$ (left boundary)};

\end{tikzpicture}
\caption[Interior and boundary finite-difference stencils]{An interior
sample such as $x_{3}$ has neighbours on both sides and uses the
symmetric central stencil (blue), whose two weights $\pm 1/(2h)$ cancel
the even-order Taylor terms and leave second-order accuracy. The first
sample $x_{0}$ has no left neighbour, so a symmetric stencil cannot be
formed; the second-order one-sided stencil (red) reads three points to
the right with weights $-3, +4, -1$ over $2h$, recovering second-order
accuracy without a left neighbour. A reader who silently drops to the
plain forward difference at the boundary loses one order of accuracy
exactly where an edge effect is already most likely to bite.}
\label{fig:vol02ch05:onesided-stencil}
\end{figure}
